\chapter{Database and audit log}\label{ch:db}

\section{Storage engine}
The audit log is a single SQLite file~\cite{sqlite}, \fpath{backend/data/audit.db}, or a file in the directory named by \code{SDCV_DATA_DIR}. It is created with \code{CREATE TABLE IF NOT EXISTS} on first start; there is no migration tool and no ORM. \fpath{audit.py} uses the standard-library \code{sqlite3} module with one connection per call and a lock around writes.

\section{Schema}
There is one table, \code{runs} (Figure~\ref{fig:er} and Table~\ref{tab:runs}). The \code{record} column holds a JSON document that contains the full \code{RunResult} and the tool versions.

\begin{figure}[H]
\centering
% Table diagram of the audit database: one table and the JSON structure of its record column.
{\hyphenpenalty=10000\exhyphenpenalty=10000
\begin{tikzpicture}[x=1cm,y=1cm,
  col/.style={font=\footnotesize\ttfamily,anchor=west,inner sep=2pt},
  typ/.style={font=\scriptsize,anchor=west,inner sep=2pt,text=neuline}]
  % table runs
  \node[draw=neuline,fill=neufill,minimum width=6.9cm,minimum height=0.8cm,anchor=north west,font=\bfseries\ttfamily] (h) at (0,0) {runs};
  \node[draw=neuline,minimum width=6.9cm,minimum height=6.7cm,anchor=north west] (body) at (h.south west) {};
  \foreach \cn/\ct/\cy in {id/TEXT PRIMARY KEY/1, created\_at/TEXT NOT NULL/2, status/TEXT NOT NULL/3, mode/TEXT NOT NULL/4,
                        summary/TEXT/5, record/TEXT NOT NULL/6, reviewer/TEXT/7, reviewed\_at/TEXT/8, review\_note/TEXT/9}
  {
    \node[col] at ($(h.south west)+(0.15,-0.7*\cy+0.25)$) {\cn};
    \node[typ] at ($(h.south west)+(3.2,-0.7*\cy+0.25)$) {\ct};
  }
  \node[font=\scriptsize,anchor=north west,text=neuline,align=left] at ($(body.south west)+(0.1,-0.15)$)
    {indexes: \texttt{idx\_runs\_created} (\texttt{created\_at} DESC), \texttt{idx\_runs\_status} (\texttt{status})\\ no foreign keys; \texttt{INSERT OR REPLACE} on save};
  % record structure
  \node[draw=detline,fill=detfill,minimum width=6.5cm,minimum height=0.8cm,anchor=north west,font=\bfseries\ttfamily] (rh) at (8.0,0) {record (JSON)};
  \node[draw=detline,minimum width=6.5cm,minimum height=6.7cm,anchor=north west] (rb) at (rh.south west) {};
  \node[font=\footnotesize,anchor=north west,align=left,text width=6.2cm] at ($(rh.south west)+(0.1,-0.15)$)
    {\texttt{result}: the \texttt{RunResult} (status, inputs, design, checks, iterations, report, narrative, trace, llm\_attempts)\\[3pt]
     \texttt{versions}: app, python, pydantic, code\_edition, rules\_fingerprint\\[3pt]
     \texttt{llm\_model}\\[3pt]
     structured runs: \texttt{raw\_input}\\[3pt]
     chat runs: redacted \texttt{messages}, \texttt{attachment\_chars}, \texttt{image\_count}, \texttt{intake}};
  \draw[flow] ($(h.south east)+(0,-3.95)$) -- node[lbl,above] {holds JSON} ($(rh.south west)+(0,-3.95)$);
  % reviewer bracket
  \draw[decorate,decoration={brace,amplitude=4pt,mirror},color=outline] ($(h.south west)+(-0.2,-4.7)$) -- ($(h.south west)+(-0.2,-6.25)$);
  \node[font=\scriptsize,rotate=90,anchor=south] at ($(h.south west)+(-0.45,-5.5)$) {set by review endpoint};
\end{tikzpicture}}

\caption{The \code{runs} table and the structure of its \code{record} JSON document. There are no foreign keys.}
\label{fig:er}
\end{figure}

\begin{table}[H]
\centering
\caption{Columns of the \code{runs} table.}\label{tab:runs}
\begin{tabular}{@{}llL{8.2cm}@{}}
\toprule Column & Type & Content \\ \midrule
\code{id} & TEXT, primary key & 12-character hexadecimal run identifier \\
\code{created_at} & TEXT, not null & UTC ISO 8601 timestamp to the second \\
\code{status} & TEXT, not null & One of the result statuses \\
\code{mode} & TEXT, not null & \code{chat} or \code{structured} \\
\code{summary} & TEXT & One-line description, for example support, span, grades, section and outcome \\
\code{record} & TEXT, not null & JSON: \code{result}, \code{versions}, \code{llm_model} and mode-specific fields \\
\code{reviewer} & TEXT & Reviewer name from the review endpoint \\
\code{reviewed_at} & TEXT & UTC timestamp of the review entry \\
\code{review_note} & TEXT & Optional note \\
\bottomrule
\end{tabular}
\end{table}

\paragraph{Indexes and constraints.} The primary key is \code{id} and a save uses \code{INSERT OR REPLACE}. Indexes are \code{idx_runs_created} on \code{created_at} descending and \code{idx_runs_status} on \code{status}. There are no foreign keys. Statistics use \code{json_extract} over \code{record}, which scans the table.

\paragraph{Contents of \code{record}.} \code{result} is the \code{RunResult} (including \code{trace} and \code{llm_attempts}); \code{versions} holds \code{app}, \code{python}, \code{pydantic}, \code{code_edition} and \code{rules_fingerprint}; \code{llm_model} names the model or models used. A structured run adds \code{raw_input}. A chat run adds the redacted \code{messages}, \code{attachment_chars}, \code{image_count} and \code{intake}.

\section{Queries and exports}
\begin{itemize}
  \item \textbf{Search.} Runs are listed newest first with optional filters on status, mode, a text query (matched with an escaped \code{LIKE} against identifier, summary and reviewer) and reviewed state, with a limit and offset and the total count.
  \item \textbf{Statistics.} Counts by status, mode, reviewed runs, runs in the last 24 hours, runs per model and runs per narrative source.
  \item \textbf{CSV export.} The filtered list as a download, up to 5000 rows, with spreadsheet formula-injection protection.
  \item \textbf{Markdown report.} A report rebuilt from the stored, already verified result, with the code edition and rule fingerprint.
\end{itemize}

\section{Reviewer log}
\code{POST /api/runs/{run_id}/review} sets \code{reviewer}, \code{reviewed_at} and \code{review_note} for an existing run. It records a log entry only. The endpoint does not verify identity, there are no user accounts, and the entry is not a certification; the stamp on every output remains.
